% Typeset from the audited Markdown manuscript.
% Responsible author: Carptopus.
\documentclass[11pt]{article}
\usepackage[a4paper,margin=27mm]{geometry}
\usepackage[T1]{fontenc}
\usepackage[utf8]{inputenc}
\usepackage{lmodern}
\usepackage{microtype}
\usepackage{amsmath,amssymb,amsthm,mathtools}
\usepackage{needspace}
\usepackage{xcolor}
\usepackage{fancyvrb}
\usepackage[colorlinks=true,linkcolor=blue!55!black,citecolor=blue!55!black,urlcolor=blue!65!black]{hyperref}
\usepackage{fancyhdr}

\newtheorem{theorem}{Theorem}[section]
\newtheorem{lemma}[theorem]{Lemma}
\numberwithin{equation}{section}

\pagestyle{fancy}
\fancyhf{}
\fancyhead[L]{Carptopus}
\fancyhead[R]{Cyclotomic minimum degree for all orders}
\fancyfoot[C]{\thepage}
\setlength{\headheight}{14pt}
\setlength{\parindent}{0pt}
\setlength{\parskip}{0.45em}
\setlength{\emergencystretch}{4em}
\urlstyle{same}

\title{The minimum degree of nonnegative multiples\\
of cyclotomic polynomials}
\author{Carptopus\\\href{mailto:carptopus@163.com}{carptopus@163.com}}
\date{28 August 2026}

\hypersetup{
  pdftitle={The minimum degree of nonnegative multiples of cyclotomic polynomials},
  pdfauthor={Carptopus},
  pdfsubject={An all-orders solution of Steinberger's cyclotomic minimum-degree conjecture},
  pdfkeywords={cyclotomic polynomial, nonnegative coefficients, minimum degree, roots of unity, trigonometric kernel, separation certificate}
}

\begin{document}
\maketitle

\begin{center}
\fcolorbox{black!35}{black!4}{\parbox{0.88\textwidth}{\centering\small
Internal candidate manuscript, version 0.1-beta. The proof and complete
manuscript have passed project-internal zero-trust audits. External mathematical
review and formal peer review are pending.}}
\end{center}

\begin{abstract}
Let $\Phi_n(x)$ be the $n$-th cyclotomic polynomial, let $p$ be the smallest
prime divisor of $n>1$, and put $P=n/p$. Steinberger formulated the conjecture
that the lowest-degree monic polynomial with nonnegative real coefficients
divisible by $\Phi_n$ is
\[
R_{n,p}(x)=1+x^P+\cdots+x^{(p-1)P}.
\]
We prove this conjecture for every $n>1$. More precisely, every nonzero
polynomial $f\in\mathbb R_{\geq0}[x]$ divisible by $\Phi_n$ has degree at
least $(p-1)P$, and equality holds exactly for the positive scalar multiples
of $R_{n,p}$.

The proof uses an explicit trigonometric separation certificate involving only
the primitive frequencies $1,\ldots,p-1$. With $\alpha=\pi/p$, define
\[
K_p(\theta)=\sum_{k=1}^{p-1}\sin(k\alpha)\cos(k\theta).
\]
An elementary closed form shows that the shifted values
$K_p(2\pi j/n+\alpha)$ vanish at the regular $p$-gon positions and have one
strict sign at every other exponent below the proposed degree. Evaluation at
primitive $n$-th roots makes this vector orthogonal to every $\Phi_n$-multiple
of degree less than $n$. A separate unit-group reduction argument gives
uniqueness on the boundary. The argument is uniform: it does not require $n$
to be squarefree and includes even integers, primes, prime powers, and integers
with arbitrarily many prime factors.
\end{abstract}

\section{Introduction}

For an integer $n>1$, define
\begin{equation}
\mu(n)=\min\{\deg f:0\ne f\in\mathbb R_{\geq0}[x],\ \Phi_n(x)\mid f(x)\}.
\label{eq:mu}
\end{equation}
If $p$ is the smallest prime divisor of $n$ and $P=n/p$, then
\begin{equation}
R_{n,p}(x)=1+x^P+\cdots+x^{(p-1)P}=\Phi_p(x^P)
\label{eq:regular}
\end{equation}
vanishes at every primitive $n$-th root of unity. Hence $\Phi_n\mid R_{n,p}$,
and
\begin{equation}
\mu(n)\leq(p-1)P.
\label{eq:upper}
\end{equation}

Steinberger \cite[Conjecture~1]{steinberger2012} conjectured equality in
\eqref{eq:upper}, with $R_{n,p}$ as the unique monic minimizer. He proved the
conjecture when $n$ is even, when $n$ is a prime power, and, writing the other
distinct prime divisors as $q_1,\ldots,q_k$, whenever
\begin{equation}
\frac2p>\frac1{q_1}+\cdots+\frac1{q_k}.
\label{eq:reciprocal}
\end{equation}
The same paper explicitly expressed disbelief in the full conjecture and
anticipated counterexamples for orders with many large prime factors. Its
general method reduces the squarefree case to a finite separation problem in a
high-dimensional zero-sum array space.

The purpose of this paper is to give a uniform affirmative solution. The
certificate below is not obtained by solving those high-dimensional systems.
It lies in a $(p-1)$-frequency subspace and is available in closed form for
every $n$.

\Needspace{0.24\textheight}
\begin{theorem}\label{thm:main}
Let $n>1$, let $p$ be its smallest prime divisor, and put $P=n/p$. If
$0\ne f\in\mathbb R_{\geq0}[x]$ and $\Phi_n\mid f$, then
\begin{equation}
\deg f\geq(p-1)P.
\label{eq:lower}
\end{equation}
Equality holds if and only if
\begin{equation}
f(x)=cR_{n,p}(x)
\label{eq:equality}
\end{equation}
for some $c>0$. Consequently $R_{n,p}$ is the unique monic minimizer.
\end{theorem}

The proof is self-contained. Section~\ref{sec:kernel} establishes the
trigonometric kernel identity and its exact sign pattern.
Section~\ref{sec:certificate} turns the kernel into a separation certificate
for every order. Section~\ref{sec:rigidity} proves the boundary rigidity needed
for uniqueness, including nonsquarefree orders. Section~\ref{sec:proof} proves
Theorem~\ref{thm:main}. Section~\ref{sec:checks} records exact destructive
checks, and Section~\ref{sec:scope} positions the result relative to earlier
work.

\section{A sector kernel}\label{sec:kernel}

Fix a prime $p$, and put
\begin{equation}
\alpha=\frac\pi p,
\qquad
K_p(\theta)=\sum_{k=1}^{p-1}\sin(k\alpha)\cos(k\theta).
\label{eq:kernel}
\end{equation}
Angles are understood modulo $2\pi$. The open cyclic sector
$(-\alpha,\alpha)$ means the arc represented by real angles strictly between
$-\alpha$ and $\alpha$.

\begin{lemma}\label{lem:kernel}
If $\theta\not\equiv\pm\alpha\pmod{2\pi}$, then
\begin{equation}
K_p(\theta)=
\frac{\sin\alpha\,[1+\cos(p\theta)]}
     {2(\cos\theta-\cos\alpha)}.
\label{eq:closed}
\end{equation}
The zeros of $K_p$ are exactly the odd multiples of $\alpha$. Moreover,
$K_p(\theta)>0$ for $\theta\in(-\alpha,\alpha)$, and $K_p(\theta)<0$ at
every point outside this sector that is not an odd multiple of $\alpha$ modulo
$2\pi$.
\end{lemma}

\begin{proof}
Set
\[
S(z)=\sum_{k=1}^{p-1}\sin(k\alpha)z^k.
\]
The recurrence
\[
\sin((k+1)\alpha)-2\cos\alpha\sin(k\alpha)
+\sin((k-1)\alpha)=0
\]
together with $\sin(p\alpha)=0$ and
$\sin((p-1)\alpha)=\sin\alpha$ gives
\begin{equation}
(1-2z\cos\alpha+z^2)S(z)=\sin\alpha\,(z+z^{p+1}).
\label{eq:generating}
\end{equation}
Substituting $z=e^{i\theta}$, dividing away from
$\theta\equiv\pm\alpha$, and taking real parts yields \eqref{eq:closed}.
At the two removable points, \eqref{eq:kernel} gives
\begin{equation}
K_p(\pm\alpha)=\frac12\sum_{k=1}^{p-1}\sin(2k\alpha)=0.
\label{eq:removable}
\end{equation}
Away from these points, the numerator in \eqref{eq:closed} is zero exactly
when $p\theta\equiv\pi\pmod{2\pi}$, hence exactly at the remaining odd
multiples of $\alpha$. At all other points the numerator is positive. The
denominator has the sign of $\cos\theta-\cos\alpha$, which is positive
exactly in $(-\alpha,\alpha)$. This proves the zero set and sign statement.
For $p=2$, the same calculation reduces to $K_2(\theta)=\cos\theta$.
\end{proof}

\section{A uniform separation certificate}\label{sec:certificate}

Return to an arbitrary integer $n>1$. Let $p$ be its smallest prime divisor,
put $P=n/p$, and define
\begin{equation}
d=(p-1)P,
\qquad
H=\{0,P,2P,\ldots,(p-1)P\}.
\label{eq:anchors}
\end{equation}
For $0\leq j<n$, set
\begin{equation}
w_j=K_p\!\left(\frac{2\pi j}{n}+\alpha\right).
\label{eq:weights}
\end{equation}

\begin{lemma}\label{lem:sign}
The vector $w=(w_0,\ldots,w_{n-1})$ satisfies
\begin{equation}
w_j=0\quad(j\in H),
\qquad
w_j<0\quad(0\leq j<d,\ j\notin H).
\label{eq:sign}
\end{equation}
\end{lemma}

\begin{proof}
If $j=hP$, then the angle in \eqref{eq:weights} is $(2h+1)\alpha$, so
Lemma~\ref{lem:kernel} gives $w_j=0$. Conversely, an odd-multiple zero in
\eqref{eq:weights} implies $j\equiv hP\pmod n$, and the representatives in
$[0,n)$ are precisely the elements of $H$.

The angle in \eqref{eq:weights} lies in the positive cyclic sector precisely
when
\[
1-\frac1p<\frac jn<1,
\]
or equivalently $d<j<n$. Thus every nonanchor with $0\leq j<d$ lies in the
strict negative region of Lemma~\ref{lem:kernel}. If $P=1$, the
strict-negative index set is empty. If $p=2$, \eqref{eq:weights} becomes
$w_j=-\sin(2\pi j/n)$, giving the same conclusion directly.
\end{proof}

\begin{lemma}\label{lem:orthogonal}
Let $f(x)=\sum_{j=0}^{n-1}a_jx^j\in\mathbb R[x]$ and suppose
$\Phi_n\mid f$. Then
\begin{equation}
\sum_{j=0}^{n-1}a_jw_j=0.
\label{eq:orthogonal}
\end{equation}
\end{lemma}

\begin{proof}
Let $\zeta_n=e^{2\pi i/n}$. For every integer $k$ with $1\leq k<p$,
every prime divisor of $k$ is smaller than $p$ and therefore cannot divide
$n$. Hence $(k,n)=1$, so $\zeta_n^k$ is a primitive $n$-th root and
\begin{equation}
f(\zeta_n^k)=0.
\label{eq:root}
\end{equation}
Using \eqref{eq:kernel} and \eqref{eq:weights}, we obtain
\begin{align}
\sum_{j=0}^{n-1}a_jw_j
&=\sum_{k=1}^{p-1}\sin(k\alpha)
  \sum_{j=0}^{n-1}a_j
  \cos\!\left(\frac{2\pi kj}{n}+k\alpha\right)\notag\\
&=\sum_{k=1}^{p-1}\sin(k\alpha)\,
  \mathop{\rm Re}\!\left(e^{ik\alpha}f(\zeta_n^k)\right)=0.
\label{eq:fourier}
\end{align}
No squarefreeness or Chinese-remainder decomposition is used.
\end{proof}

\section{Boundary rigidity for arbitrary orders}\label{sec:rigidity}

The separation certificate will force a minimizer onto $H$. We next show that
divisibility then forces all anchor coefficients to be equal.

\begin{lemma}\label{lem:surjective}
The reduction homomorphism
\begin{equation}
(\mathbb Z/n\mathbb Z)^*\longrightarrow(\mathbb Z/p\mathbb Z)^*
\label{eq:reduction}
\end{equation}
is surjective.
\end{lemma}

\begin{proof}
Write $n=p^em$, where $e\geq1$ and $(p,m)=1$. Given
$b\in(\mathbb Z/p\mathbb Z)^*$, choose a unit $b'$ modulo $p^e$ reducing to
$b$. If $m>1$, the Chinese remainder theorem gives an integer $a$ satisfying
\[
a\equiv b'\pmod{p^e},
\qquad
a\equiv1\pmod m.
\]
Then $(a,n)=1$ and $a\equiv b\pmod p$. The case $m=1$ is immediate.
\end{proof}

\begin{lemma}\label{lem:rigidity}
Suppose $0\ne f\in\mathbb R_{\geq0}[x]$, $\Phi_n\mid f$,
$\deg f\leq d$, and $\operatorname{supp}(f)\subseteq H$. Then
\begin{equation}
f(x)=cR_{n,p}(x)
\label{eq:rigidity}
\end{equation}
for some $c>0$.
\end{lemma}

\begin{proof}
Write $f(x)=g(x^P)$, where $\deg g\leq p-1$. If $\zeta_n$ is a primitive
$n$-th root, then $\zeta_p=\zeta_n^P$ has order $p$. For every
$a\in(\mathbb Z/n\mathbb Z)^*$,
\begin{equation}
0=f(\zeta_n^a)=g(\zeta_p^a).
\label{eq:groot}
\end{equation}
By Lemma~\ref{lem:surjective}, the exponents $a\bmod p$ range over all
nonzero residue classes. Therefore $g$ vanishes at every primitive $p$-th
root. Hence $\Phi_p\mid g$. Since $g\ne0$,
$\deg g\leq p-1=\deg\Phi_p$, and
$\Phi_p=1+x+\cdots+x^{p-1}$, we have $g=c\Phi_p$. The coefficient
conditions force $c>0$, proving \eqref{eq:rigidity}.
\end{proof}

\section{Proof of the main theorem}\label{sec:proof}

The polynomial $R_{n,p}=\Phi_p(x^P)$ vanishes at every primitive $n$-th
root: if $\zeta_n$ is primitive, then $\zeta_n^P$ is a nontrivial $p$-th
root. Thus $\Phi_n\mid R_{n,p}$, and \eqref{eq:upper} holds.

Conversely, suppose that $0\ne f\in\mathbb R_{\geq0}[x]$,
$\Phi_n\mid f$, and $\deg f\leq d$. Because $d<n$, write
\[
f(x)=\sum_{j=0}^{n-1}a_jx^j
\]
by padding with zero coefficients. Lemma~\ref{lem:orthogonal} gives
\begin{equation}
0=\sum_{j=0}^{d}a_jw_j.
\label{eq:separate}
\end{equation}
Every term in \eqref{eq:separate} is nonpositive by Lemma~\ref{lem:sign}.
At each nonanchor $0\leq j<d$, it is strictly negative if $a_j>0$.
Therefore all such coefficients vanish and
$\operatorname{supp}(f)\subseteq H$. Lemma~\ref{lem:rigidity} now gives
$f=cR_{n,p}$ with $c>0$.

It follows simultaneously that no nonzero admissible polynomial has degree
less than $d$, and that every degree-$d$ admissible polynomial is a positive
scalar multiple of $R_{n,p}$. This proves Theorem~\ref{thm:main}.

\section{Exact destructive checks}\label{sec:checks}

The proof above is symbolic and does not depend on computation. We nevertheless
used exact and high-precision checks to attack the most likely failure modes.
The accompanying verifier
\texttt{verification/verify\_full\_n\_sector\_certificate.py} checks:
\begin{itemize}
\item every $2\leq n\leq1000$, separated into primes, prime powers, mixed
  squarefree orders, and mixed nonsquarefree orders;
\item every anchor by direct high-precision finite summation for
  $2\leq n\leq300$, selected anchors above that range, and every nonanchor
  prefix and tail sign for $2\leq n\leq1000$;
\item exact gcd checks for primitivity of all frequencies $1,\ldots,p-1$;
\item surjectivity of
  $(\mathbb Z/n\mathbb Z)^*\to(\mathbb Z/p\mathbb Z)^*$;
\item for every $2\leq n\leq180$, divisibility of $R_{n,p}$ and 6410
  direct orthogonality checks on shifted $\Phi_n$-basis elements;
\item direct finite-sum versus closed-form evaluations at 24 nontrivial angles;
\item a phase-removal negative control and a wrong-prime negative control.
\end{itemize}

Run it from this entry directory with
\begin{Verbatim}[fontsize=\footnotesize]
python -X utf8 .\verification\verify_full_n_sector_certificate.py
\end{Verbatim}
The verifier SHA-256 is
\begin{center}
\small\texttt{62E512FD23DB231D05E08BDB074F9F36C37CC5FE0AC7337CA88BB0C249895AB3}.
\end{center}
All checks pass. Removing the phase destroys the anchor zeros, while replacing
the smallest prime by an arbitrary prime can introduce nonprimitive low
frequencies, as expected. These computations are destructive controls only;
the all-orders conclusion rests on Lemmas~\ref{lem:kernel}--\ref{lem:rigidity}.

\section{Relation to earlier work and scope}\label{sec:scope}

Steinberger \cite{steinberger2012} introduced problem \eqref{eq:mu},
formulated Conjecture~1, proved the reciprocal sufficient condition
\eqref{eq:reciprocal}, and developed a squarefree array separation criterion.
The present proof uses the same broad dual-separation idea, but its certificate
is the explicit low-frequency vector \eqref{eq:weights}, whose membership and
sign are established directly and uniformly for every order.

Reddy \cite{reddy2011} studies a different restricted problem involving
$0,1$-polynomials and minimal vanishing sums, principally for even orders with
few distinct prime factors. Kepka and Korbel\'a\v{r} \cite{kepka2012} give degree
bounds for nonnegative multiples of a general polynomial and solve their cubic
case. {\L}aba and Zakharov \cite{laba2025} use Steinberger's bound in the study of
integer tilings. None of these results states Theorem~\ref{thm:main}.

Two earlier public-beta notes by the present author treated portions of the
same problem. The first \cite{carptopus-fixed} gave exact dense certificates
for the first two orders outside Steinberger's theorem. The second
\cite{carptopus-tail} constructed symbolic tail-projection certificates for
infinite prime families outside \eqref{eq:reciprocal}. Theorem~\ref{thm:main}
strictly subsumes the mathematical conclusions of both notes; their
certificate constructions remain separate computational and structural
records, but they are not needed for the proof here.

A targeted literature refresh on 28 August 2026 checked the original article,
its current OpenAlex forward-citation set, arXiv, and exact-phrase and
method-level searches. OpenAlex listed ten citing works; the newest was
\cite{laba2025}, and the only directly adjacent lowest-degree papers were
\cite{reddy2011,kepka2012}. We found no public source claiming an all-orders
proof, the sector certificate \eqref{eq:weights}, or an affirmative resolution
of Conjecture~1. This is a bounded public search, not proof of global priority:
unindexed, unpublished, or concurrent work could still affect the novelty
assessment.

\section{AI-assisted research and writing disclosure}

OpenAI Codex was used for computational exploration, proof development and
organization, adversarial checking, literature search, and language editing.
Carptopus selected and reviewed the retained mathematical claims and assumes
responsibility for the manuscript.

\begingroup
\small
\begin{thebibliography}{9}

\bibitem{steinberger2012}
J. P. Steinberger,
``The lowest-degree polynomial with nonnegative coefficients divisible by the
$n$-th cyclotomic polynomial,''
\emph{Electronic Journal of Combinatorics} \textbf{19}(4) (2012), Paper P1.
\href{https://doi.org/10.37236/2755}{doi:10.37236/2755}.

\bibitem{reddy2011}
A. S. Reddy,
``The lowest degree $0,1$-polynomial divisible by cyclotomic polynomial,''
arXiv:1106.1271v2 (2011).
\url{https://arxiv.org/abs/1106.1271}.

\bibitem{kepka2012}
T. Kepka and M. Korbel\'a\v{r},
``The lowest-degree polynomials with non-negative coefficients,''
arXiv:1210.6868 (2012).
\url{https://arxiv.org/abs/1210.6868}.

\bibitem{laba2025}
I. {\L}aba and D. Zakharov,
``On the minimal period of integer tilings,''
\emph{Bulletin of the London Mathematical Society} \textbf{57}(4) (2025),
1160--1170.
\href{https://doi.org/10.1112/blms.70023}{doi:10.1112/blms.70023}.

\bibitem{carptopus-fixed}
Carptopus,
``Exact certificates for the first two cases not covered by Steinberger's
cyclotomic minimum-degree theorem,'' Zenodo, v0.1-beta (2026).
\href{https://doi.org/10.5281/zenodo.22124096}{doi:10.5281/zenodo.22124096}.

\bibitem{carptopus-tail}
Carptopus,
``Tail-projection certificates for cyclotomic minimum-degree polynomials beyond
the reciprocal condition,'' Zenodo, v0.1-beta (2026).
\href{https://doi.org/10.5281/zenodo.22127365}{doi:10.5281/zenodo.22127365}.

\end{thebibliography}
\endgroup

\end{document}
